\section{Introduction}
\label{sec:introduction}

Autonomous-driving perception requires an understanding of the objects
and surfaces surrounding a vehicle. \Gls{lidar} provides
three-dimensional measurements of the environment that support tasks
such as \gls{object-detection} and \gls{semantic-segmentation}
\parencite{geigerAreWeReady2012a,behleySemanticKITTIDatasetSemantic2019}.
Developing learning-based methods for these tasks requires annotated
data, but collecting and labelling automotive \glspl{point-cloud} is
intensive. Furthermore, uncommon traffic situations are difficult
to capture in sufficient numbers through ordinary data collection.
These constraints motivate the generation of \gls{synthetic-data},
which offers opportunities to produce annotated observations and
construct controlled scenarios
\parencite{hurlPreciseSyntheticImage2019,
	manivasagamLiDARsimRealisticLiDAR2020,xiaoTransferLearningSynthetic2022}.

The potential of synthetic data depends on how the data is
generated. Virtual environments can provide scene control and automatic
annotations, as demonstrated by PreSIL and SynLiDAR
\parencite{hurlPreciseSyntheticImage2019,
	xiaoTransferLearningSynthetic2022}.
Other approaches draw more directly on real observations. LiDARsim
combines reconstructed scenes with physical simulation and a learned
model of missing returns, while LiDARGen learns to generate scans from
examples
\parencite{manivasagamLiDARsimRealisticLiDAR2020,
	zyrianovLearningGenerateRealistic2022}.
Together, these approaches show different balances between
explicit control over the scene and dependence on data observed in the real world.
Understanding these differences provides a basis for assessing which
generation approaches can meet the requirements of a particular
perception task.

For automotive LiDAR, those requirements extend beyond the geometry of
the scene. A scan samples visible surfaces from a particular viewpoint,
and its structure depends on sensor sampling, \gls{occlusion} and
missing returns. LiDARsim illustrates the importance of modelling this
observation process, while SynLiDAR identifies differences in appearance
and sparsity as components of the \gls{domain-gap} between synthetic
and real point clouds
\parencite{manivasagamLiDARsimRealisticLiDAR2020,
	xiaoTransferLearningSynthetic2022}.
Consequently, the assessment of synthetic data must connect the
properties of generated scans with their intended use. LiDARGen
evaluates scan generation and completion, whereas PreSIL and SynLiDAR
also investigate the contribution of synthetic data to training
perception models
\parencite{zyrianovLearningGenerateRealistic2022,
	hurlPreciseSyntheticImage2019,xiaoTransferLearningSynthetic2022}.
These evaluations provide different forms of evidence, making it
necessary to distinguish similarity to real world observations from
demonstrated benefits for learning on real world data.

This narrative review examines that connection through selected
literature on synthetic-data generation and its application to
scene level automotive LiDAR perception. General generation approaches
provide the conceptual foundation, while LiDAR-specific studies form
the basis for discussing object detection and semantic segmentation.
The scope includes virtual-scene simulation, reconstruction-based
simulation and learned scan generation. Scan completion is considered
where it clarifies the capabilities of a method, while its results are
distinguished from evidence obtained by training with newly generated
data. This scope supports a comparison of generation capabilities,
sensor-specific requirements and evaluation practices without assuming
that results from different tasks and datasets are directly comparable.

The review first examines the principal generation approaches and
their implications for control, coverage and dependence on real data.
It then considers their application to automotive LiDAR, focusing on
scene construction, sensor modelling and annotation consistency.
These requirements provide the basis for examining evaluation metrics
and evidence of transfer to real perception tasks. The discussion
subsequently brings these findings together to identify limitations
and directions for further investigation. The review therefore begins
with the generation approaches that underpin the construction and
assessment of synthetic LiDAR data.
